\chapter{Phase 1: Foundations (The Stateless Phase)}

In this phase, we focus on understanding how to control Large Language Model (LLM) outputs deterministically. This is the bedrock of Harness Engineering: moving away from ``prompt and pray'' towards rigorous software control.

\section{Step 1: API Mastery}
The goal of this step is to understand core LLM APIs and parameters. A Harness Engineer must own the model's behavior through deterministic settings.

\begin{lstlisting}[caption={step1\_api\_mastery.py}, label={lst:step1}]
"""
PHASE 1 - STEP 1: API Mastery
Harness Pillar: Cognitive Framework
Goal: Understand how to deterministically control Claude's output using API parameters.
"""

import anthropic
from dotenv import load_dotenv

load_dotenv()
client = anthropic.Anthropic()

# 1a. Basic call - observe the raw structure
response = client.messages.create(
    model="claude-opus-4-5-20251101",
    max_tokens=256,
    messages=[{"role": "user", "content": "What is 17 x 24?"}],
)
print(f"Answer: {response.content[0].text}")

# 1b. Temperature = 0 (deterministic output)
for _ in range(3):
    r = client.messages.create(
        model="claude-opus-4-5-20251101",
        max_tokens=64,
        temperature=0,
        messages=[{"role": "user", "content": "Give me one word for 'happy'."}],
    )
    print(r.content[0].text.strip())

# 1c. Stop sequences - force early termination
r = client.messages.create(
    model="claude-opus-4-5-20251101",
    max_tokens=256,
    stop_sequences=["STOP"],
    messages=[{"role": "user", "content": "Count from 1 to 10. After each number write 'STOP' on a new line."}],
)
print(r.content[0].text)

# 1d. System prompt - define the Cognitive Framework
SYSTEM = """You are a financial data analyst assistant.
Rules:
1. ONLY answer questions about financial data.
2. ALWAYS cite your reasoning step-by-step.
3. Format all monetary values as: $X,XXX.XX
"""
r = client.messages.create(
    model="claude-opus-4-5-20251101",
    max_tokens=256,
    system=SYSTEM,
    messages=[{"role": "user", "content": "What's the PE ratio of Apple?"}],
)
print(r.content[0].text)
\end{lstlisting}

\subsection*{Key Takeaways}
\begin{itemize}
    \item Always set \texttt{temperature=0} for deterministic agent actions.
    \item Use \texttt{stop\_sequences} to force the model to produce parseable output.
    \item The system prompt IS your agent's Cognitive Framework---define its identity and boundaries here.
\end{itemize}

\section{Step 2: Structured Outputs}
Harness Engineers never parse free text with regular expressions if they can avoid it. Instead, they contract the model to emit machine-readable structured data using schemas like Pydantic.

\begin{lstlisting}[caption={step2\_structured\_outputs.py}, label={lst:step2}]
"""
PHASE 1 - STEP 2: Structured Outputs
Harness Pillar: ACI (Agent-Computer Interface)
Goal: Force Claude to reply in strictly validated JSON schemas using Pydantic.
"""

from typing import Literal
import anthropic
from pydantic import BaseModel, Field
from dotenv import load_dotenv

load_dotenv()
client = anthropic.Anthropic()

class Company(BaseModel):
    name: str
    founded_year: int
    headquarters: str
    main_product: str

response = client.messages.parse(
    model="claude-opus-4-5-20251101",
    max_tokens=256,
    output_format=Company,
    messages=[{"role": "user", "content": "Tell me about Anthropic, the AI safety company."}],
)

company: Company = response.content[0].parsed
print(f"Name: {company.name}, Founded: {company.founded_year}")

# Decision schema - agent decision-making
class AgentDecision(BaseModel):
    action: Literal["approve", "reject", "escalate"]
    reason: str
    confidence: float = Field(ge=0.0, le=1.0)

SCENARIO = "User wants to withdraw $50,000. Balance is $52,000. Account is 2 years old."

response = client.messages.parse(
    model="claude-opus-4-5-20251101",
    max_tokens=256,
    system="You are a bank fraud detection agent.",
    output_format=AgentDecision,
    messages=[{"role": "user", "content": f"Evaluate: {SCENARIO}"}],
)

decision: AgentDecision = response.content[0].parsed
print(f"Action: {decision.action}, Confidence: {decision.confidence:.0%}")
\end{lstlisting}

\subsection*{Key Takeaways}
\begin{itemize}
    \item \texttt{client.messages.parse()} is the cleanest pattern for structured outputs.
    \item Use \texttt{Literal} for constrained choices (enums).
    \item Structured outputs are the foundation of reliable agentic pipelines.
\end{itemize}

\section{Step 3: The ReAct Loop}
The ReAct (Reason + Act) loop is the heart of every AI agent. It is a cycle where the model thinks, acts (calls a tool), observes the result, and repeats until the task is complete.

\begin{lstlisting}[caption={step3\_react\_loop.py}, label={lst:step3}]
"""
PHASE 1 - STEP 3: The ReAct Loop (Reason + Act)
Harness Pillar: Standard Workflows & Refinement Loops
Goal: Build the core agentic loop from scratch without frameworks.
"""

import math
import json
import anthropic
from dotenv import load_dotenv

load_dotenv()
client = anthropic.Anthropic()

# Tool definition (The ACI)
TOOLS = [{
    "name": "calculator",
    "description": "Perform mathematical calculations.",
    "input_schema": {
        "type": "object",
        "properties": {
            "expression": {"type": "string", "description": "Python math expression"}
        },
        "required": ["expression"]
    }
}]

def run_agent(user_query: str) -> str:
    messages = [{"role": "user", "content": user_query}]
    while True:
        response = client.messages.create(
            model="claude-opus-4-5-20251101",
            max_tokens=1024,
            tools=TOOLS,
            messages=messages,
        )
        messages.append({"role": "assistant", "content": response.content})

        if response.stop_reason == "end_turn":
            return response.content[-1].text

        if response.stop_reason == "tool_use":
            tool_results = []
            for block in response.content:
                if block.type == "tool_use":
                    # Execute tool (omitted for brevity)
                    result = "4" 
                    tool_results.append({
                        "type": "tool_result",
                        "tool_use_id": block.id,
                        "content": result,
                    })
            messages.append({"role": "user", "content": tool_results})
\end{lstlisting}

\subsection*{Key Takeaways}
\begin{itemize}
    \item The ReAct loop is a \texttt{while} loop keyed on \texttt{stop\_reason}.
    \item Always append the assistant's full content to the message history to give it ``memory.''
    \item Always cap iterations to prevent runaway loops.
\end{itemize}
